\documentclass{article}
\usepackage{graphicx} % Required for inserting images

\title{Weihnachtsgeschichte 2024}
\author{Simon Grätz}
\date{December 2024}

\begin{document}

\maketitle
\tableofcontents
\section{Deutsch}
\subsection{Vorwort.}
Es war ein mal vor langer, langer Zeit. An einem weit entfernten Ort. Zu
einer unbekannten Zeit. Also ungefähr vor ein paar Jahren 1-2 Wochen vor Weihnachten. Aus Mangel an Zeit, Geld
und Fähigkeit der deutschen Rechtschreibung entschloss ich mich meine Lieben mit einem Machwerk zu bedenken,
dessen neuste Iteration nun auch dir wertem Leser zu händen lieget. Aus einer fixen Idee wurde eine Tradition
und nun verschicke ich die frohe Botschaft nun schon seit einigen Jahren. Wie jedes Jahr gilt, schaltet euer
Hirn aus, es ist alles die reinste Wahrheit und bestimmt nicht erfunden. Oder etwa doch? Und vorallem, wer
Rechtschreibfehler findet, der möge sie zählen, quadrieren und das Ergebnis auf eine Servierte in ihrem oder
seinem Lieblingsrestaurants schreiben. Viel Spaß dabei. Nachdem das aus der Welt ist, möchte ich nun alle bitten
ihr Glas zu erheben und ja noch oben lassen, noch ein bisschen ja jetzt prost allerseits. Und jetzt wünsche ich
viel Spaß beim Lesen meines Stusses, der bestimmt nicht unter Einfluss von Bewusstseinserweiterndem Hopfentees
zustande kam. (Gott die deutsche Sprache ist was schönes, da wird die AI sich freuen das in englisch zu
übersetzen). 
\subsection{
  Ein seltsamer Brief.
}
\noindent 
Wie ihr alle natürlich wisst bin ich ein begeistertder Seefahrer. Weder die
Badewanne noch das Waschbecken sind von mir sicher. Nicht desto trotz war ich geringfügig überrascht, als eines
Tages ein Brief in meinen Briefkasten fand. Eigentlich hoffte ich auf eine Bestellung feinster Schoki, aber
leider nein,
es war nur ein Brief. Zusätzlich zu dem bitteren Geschmack von Tinte anstatts süßer Kakaonoten hatte sich der
Absender auch noch in der Person geirrt. Der Brief war auch noch auf Französisch und ihr wisst ich bin bin ein
Akademiker - soll heißen ich bilde mir ein jede Sprache ein bisschen zu verstehen. Schrecke aber vor
körperlicher Arbeit zurück wie der Teufel vor dem Weihwasser. Was die Nachricht bedeutete konnte ich zu diesem
Zeitpunkt noch nicht umreisen. Allerdings verstand ich eine Sache sehr genau. Zugtickets für eine Reise 1.
Klasse sollte man nicht verkommen lassen. Auch wenn man Name weder \textit{Adrien} noch \textit{de Gerlach} ist.
Ich konnte aus dem Ziel der Tickets schließen, dass sich beim Absender um einen Belgier und nicht um einen
Franzosen handelte, weil das Ziel eine belgische Hafenstadt war. Und da die Schokolade, die ich bestellte aus
der Schweiz stammte und es in Belgien durchaus brauchbares Bier zu geben scheint, entschied ich demokratisch die
Einladung anzunehmen. Die Abstimmung hatte insgesamt einen Teilnehmer.
\\
\noindent 
Eine Woche später erreichte ich Antwerpen. Mega interessant - ich hatte es immer Antwerben ausgesprochen. Mit
einem B. Lustig. Naja Da die Zugtickets nicht an mich sondern an eine andere Person gerichtet waren, machte ich
mich keine Hofnungen, weiter als bis hier her zu kommen. Mein Plan war es mir die Stadt und einige belgische
Bierkrüge von innen anzusehen und dann wieder zurück zu reisen. Traurige Annektodte am Rande - das Flugticket
zurück kostet natürlich einen Bruchteil von den Zügen. Wusstet ihr, dass Antwerpen die zweit größte Stadt
Belgiens ist? Zum Mittagessen gab es Muscheln mit Fritten und wenn ich geant hätte was mir bevor stand, dann
hätte ich mir eventuell noch ein 2tes Bier schmecken lassen. Ja okay ein 3tes Bier. Aber es war Urlaub. Dont
judge me! Gut gelaunt folgte ich nach dem Essen und einer kleinen Sightseeing tour (In Antwerpen steht laut
eigener Angabe das älteste Museum der Welt!), den Angaben aus dem Brief, die ich immer noch nicht komplett
übersetzt hatte. Mein Weg führte mich in den \textit{ größten Hafen Europas } (Meine Fresse was geht mit Antwerpen
up?). Vorbei an riesigen Kreuzfahrtschiffen führte mich mein Handy oder besser google maps zu einem prächtigen,
wenn auch nicht ganz ungewöhnlichem Schiff. Geschäftiges Treiben um, auf und vor dem Segelschiff lies mich
vermuten, dass das Auslaufen wie wir \textit{Seebären} sagen würden, kurz bevor stand. Interessiert und
geringfügig fehl am Platz wirkend stand ich nun da. Mit meinem Koffer und dem Brief. Belgische Seeleute riefen
sich Dinge an den Kopf. Dinge ein und verladen.
\\
\noindent 
Ich wollte gerade kehr machen, als eine altmodisch aussehender
Kapitän mit Mütze und allem mich erspäte. Ich erparre euch und mir den Wortlaut wieder zu geben, weil verstanden
habe ich ihn nicht. Jedenfalls winkte er mich an Bord, ich stolperte unbeholfen über eine schmale Holzbrücke und
stand nun vor ihm, Kapitän \textit{Georges Lecointe}. Er salutierte er, dann hielt er mir herzlich die Hand hin.
Ich schüttelte diese, man will ja nicht unhöflich sein. Und dann deutete ich auf seine uniform und meine
Touriste gadarobe und versuchte mich an einem Witz auf französisch. Der Kapitän began gröllend zu Lachen. Ob mit
oder über mich war mir nicht bewusst, jedoch zweifelte ich nun auch an seinen Französischkenntnissen. Ein
zweiter Mann stellte sich geschwind zu uns, ergriff begeistert meine nun wieder freie Hand und schüttelte mich.
In altmodischem Englisch stellte er sich als Frederik Cook, Arzt aus New York vor. Man hörte es ihm direkt an,
dass seine Immigration in die Staaten recht frisch war. Und als ich auf Deutsch fragte ob er Friedrich oder
Frederik genannt werden wollte, war er sichtlich erfreut das ich \textit{erstaunlicherweise auch} Deutsch spräche.
In dem Moment wurde mir heiß hinter den Ohren. Auch? Als nächstes kam ein sehr großer, sehr breit muskulöser
Mann zu mir. Er knurrmurmelte etwas in pseudofranzösisch. Ich blickte den Doktor an. Dieser richtete sich an
Herrn Admunson und meinte in gestückelten English ich könne auch \textit{Norwegisch}. Die Miene von Admunson
erhellte sich und er knurrmurmelte mir nun einen Fetzen hin, der erstaunlicherweise besser verständlich war als
das Französisch zuvor. Und ich spreche kein Wort Norwegisch. Im Kopf versuchte ich zu rekapitulieren. Diese
Leute dachten, ich sei ein französisch sprachiger Belgier, der auch Englisch und Norwegisch und Deutsch spricht.
Ich musste mir eine Grinsen verkneifen. Admunson blickte mich immer noch an als würde er einen Antwort erwartet,
ich hob meine Hand als ob einer Geste, mit der ich dieser Farse einhalt gebieten wollte, doch da zog der Kapitän
mich hinter sich her ins Innere des Schiffes.
\\
\noindent 
Durch einen engen Gang vorbei and Kajüten zeigte er mir ein paar räume, einen Speiseraum, eine art Karten raum
und schließlich öffnete er eine Tür schob mich hinein, verabschiedete sich und ich stand in meiner Kajüte. Das
Schiff war zwar beindruckend groß, aber trotzdem war ich überrascht wie luxuriös es aussah hier. Okay ich hatte
praktisch ein Bett, einen Tisch und ein Fenster bzw. ein Bullauge. Aber da war sogar eine weitere Tür in ein
kleines Badezimmer. Ich hätte nicht gedacht das Matrosen auf einem altmodischen Segelschiff so eine Luxuskabine
bekommen. Gesagt und auf die Zunge gebissen. Die Uniform, die eindeutig mir zu gedacht war sah alles andere als
nach einem Matrosen aus. Im Kopf überlegte ich panisch ob ich sogar noch mehr Lametta dran hatte als der, wie
ich dachte \textit{Kapitän}. Mit wackeligen Beinen zog ich mir die Verkleidung an. Welche Strafe steht auf
imitieren eines hohen Würdenträgers zu See auf einem Segelschiff? Kielholen? Ich weiß zwar nicht mehr was das
ist. Aber ich habe mir gemerkt, dass es nicht schön ist. Naja, ist ja alles ein großes Missverständnis. Der,
hoffentlich, Kapitän hat mir das Zeug ja quasi aufgezwungen. Als ich den engen Gang wieder zurück gehe und durch
eine kleine Tür ins freie trette empfängt mich eine Applaus. Na toll. Ich hebe beschwichtigend den Arm und
steuere auf den (FUUUUUUUUUUUUUCK ich hab mehr Lametta als der Dude) Kapitän zu. Ich habe eindeutig mehr goldene
Kringel und Streifen auf meiner Jacke und meine Mütze ist gefühlt doppelt so groß wie die von dem. Der Doktor
stellt mir alle nochmal förmlich vor, und reicht dem jeweiligen dabei ein Glas Champagner. Das ist also Doktor
Friedrich Koch, Roarl Admunson der 1. Offizier, 2 weitere Offizier, deren Namen ich instant vergesse und ja doch
Kapitän \textit{Georges Lecointe}. Ich weiß nicht ob ihr es wusstet, aber auf Schiffen ist es scheinbar normal,
dass der Kapitän nicht der Oberbabo ist - das bin nämlich ich der Expeditionsleiter. In diesem Moment beginnen
wir auszulaufen.
Ich blicke sehnsüchtig zurück. Wenn man die Augen ganz fest zusammen kneift kann man einen Belgischen
Expeditionsleiter sehen, der verzweifelt am Pier Hammpelman sprünge veranstalltet, weil man ihn vergessen hat.
\\
\subsection{
  Auf hoher See.
}
\subsubsubsection{Meine erste Woche auf See.}
\noindent 
Mittlerweile musste ich nicht mehr im Minutentakt kotzen. Nein wirklich man gewöhnt
sich vervorragend an das Schaukeln. Es hat sich herausgestellt, dass ich tatsächlich mit meiner Kabine dem
durchschnittlichen Matrosen weit voraus bin. Unter Deck gibt es einen stinkenden, feucht warmen Raum der an
einen Schlafsaal in einem Hostel erinnert. Darin hängen Hängematte und Seesäcke. Die Offiziere und eine Reihe
von Wissenschaftlern haben Kajüten, die etwas größer als meine Chefsuite sind, jedoch von 4 Männer gleichzeitig
bewohnt werden. Auch den Luxus eines eigenen Bads hat außer mir nur der Kapitän - ein feiner Kerl. Er plappert
mich in einer Tour auf Französisch an. Ich bilde mir langsam ein auch den einen oder anderen Brocken zu
verstehen. Genau wie Norwegisch. Ist eingentlich eine todesleichte Sprache. Ein nordischer Hühne kommt zu dir
und sagt irgendwas mit so wenig Worten wie möglich. Du nickst und fügst vielleicht noch ein Ja in einer beliegen
Sprache hinzu. Was auch immer Admunson gesagt hat, er meldet es mir nur aus Pflicht und Hierachientreue, aber
brauchen tut er mich nicht. Der könnte vermutlich das ganze Schiff alleine steueren.
\\
\noindent 
Appropo steuern. Wohin steueren wir denn gerade? Nun ich habe absolut keine Ahnung. Ich weiß, dass in keine
Himmelsrichtung Land in Sicht ist und daher könnten wir genauso auf dem Mittelmehr oder dem Indischen Ozean
sein. Wobei der Ärmelkanal oder vielleicht auch schon der Nordatlanik mehr Sinn machen würde. Merkt ihr schon
was für ein Krasser Navigator ich bin, oder?
\\
\subsubsubsection{Meine 4te Woche auf See.}
\noindent 
Langsam befürchte ich, dass ich hier nicht auf einem kurzen Vergnügungstrip bin. Außerdem habe ich auch
Gewissenbisse, dass ich sowenig zur Reise betragen kann. Klar alle behandeln mich wie einen König. Der Dude den
ich versehenlich vertrete ist wahrscheinlich der Besitzer vom Schiff und bezahlt hier alles, aber wenn der
selber mitfährt dann muss der doch auch einen Job haben, oder? Der Kapitän verbringt die meiste Zeit auf Deck
oder in seiner Kabine deswegen habe ich nun angefangen mich die halbe Nacht im Karteraum zu verschanzen. Von
außen mag es so aussehen, als würde ich Seekarten studieren oder so, aber in Wirklichkeit lese ich sämtliche
Wörterbücher und alles was auch nur im entferntestens Deutsch oder Englisch den vielen anderen Sprachen an Bord
gegenüber stellt. Es ist herrlich irronisch, wie ich versuche Französch und Norwegisch zu lernen auf einem
Schiff voller Belgier und Norweger ohne das diesen auffällt, dass ich weder die eine noch die andere Sprache
verstehe.
\\
\noindent 
Tagsüber verbringe ich dann viel Zeit auf Deck und höre den Leuten zu. Oft unterhalte ich mich auch mit Fritz
dem Arzt. Es ist angenehm mich nicht verstellen zu müssen,
den er denkt ja ich kann halt sehr gut Deutsch und so brauche ich mir keine komplizierte Geschichte überlegen.
Es gab schon Situationen, da konnte ich ihm helfen briefe an seine amerikanische Frau auf Englisch zu schreiben.
Weil ich natürlich nicht direkt Fragen kann wo wir hinfahren, erkundige ich mich indirekt bei dem Arzt bzgl.
seinem Vorratsstand. Wenn wir an Land gehen würde, dann wüsste ich wo wir sind und wie unser Kurs aussieht.
Für einen irren Moment denke ich dann, dann könnte ich unser Ziel approximieren. Als ob ich nicht bei der ersten
Gelegenheit an Land fliehen würde. \textit{Leider} sind wir sogut mit Vorräten ausgestattet, dass eine Anhalten
gar nicht nötig ist. Ich muss also weiter beobachten und lauschen und so tun als ob ich wüsste was vor sich
geht.
\\
\subsubsubsection{Meine 5te Woche auf See.}
\noindent Ach du Scheiße!!! Das waren meine Worte, als ich heute Morgen auf das Deck ging. Bzw. \textit{Oh, putain de
  merde}.
Ja mein Französisch ist langsam brauchbar genug um mich mit den Leuten zu verständigen. Lacht
mich nicht aus aber \textit{Adrien de Gerlach}, die Person die ich inpersioniere, hat zwar einen Französich
klingenden Namen, kommt aber anscheinden aus dem deutsch sprachigen Teil des Landes. Und weil alle sofort
gecheckt
haben, dass meine Französisch nicht zum Brot kaufen reicht haben sie mich halt als einen von diesen blöden
Deutsch-Belgier abgetan. Und niemand ist auf die Idee gekommen, dass das weired ist. Zurück zu meinem
ausgezeichneten Fluch auf Französisch. Eine Eisscholle treibt an uns backbord vorbei. Ich glaube ich weiß wo wir
hinsteuern, richtung Süden. Gaanz Richtung süden. Der Name des Schiffes kam mir die ganze Zeit schon bekannt
vor, aber jetzt fällt es mir wie Schuppen von den Augen. Ich bin auf der Belgica!
\\
\subsubsubsection{Meine letzte Woche auf See.}
\noindent Es war nie sehr warm an Bord. Bzw. es war schon warm in meiner Kabine, die die Abgasrohre der Öfen als
Fußbodenheizung hat. Aber direkt im Gang daneben oder auf Deck wurde es einem unmissverständlich klar, dass man
sich auf hoher See befindet. Aber die Temperaturen waren nie extrem kalt, sondern frisch, belebend. Doch seit
ein paar Tagen merkt man, dass wir uns weiter und weiter dem Ziel der Expeditions nähern. Inzwischen weiß ich
was uns bevor steht. Die Belgica is eines dieser Schiffe, dass aufgebrochen ist um den Südpol zu erreichen. Da
hab ich mir ja schön was eingebrockt. Mit einer säuerlichen Miene verlasse ich meine Kabine und gehe auf Deck.
Vorbei sind die Zeiten für die schicke Uniform. Längst hängt eine langer, warmer Mantel über meinen Schultern.
Eine dicke
Strickmütze ziert meine Haupt. Ich bin schlecht gelaunt, weil ich natürlich schlecht jetzt zugeben kann, dass
ich gar nicht \textit{Adrien de Gerlach} bin. Folgerichtig kann ich die Mission nicht einfach zum Umkehren
bewegen, auch wenn wir gerade Nase voran ins Packeis der Antarktik steuern. Und jedesmal wenn ich vorsichtig
andeute ob die Leute sich sorgen wegen dem Eis machen antworten sie nur mit einer todesverachtenden Coolheit von
Seefahreren, die schon gegen Riesenkraken den Mahlstrom und den Klabauterman gekämpft haben. Der Kapitän begrüßt
mich strahlend und wir wechseln ein paar Worte. Dann hält er mir eine Karte vor die nase und möchte meine
Meinung dazu wissen. Scherzhaft sage ich \textit{Nach Norden}? Johlendes Lachen der Besatzung ob meines
hervorrangendenden Scherzes. Ich nicke nur und schließe die Augen. Ich zähle bis 10. Ich öffne meine Augen.
Alles ist in Ordnung. Ich nicke dem Kapitän zu und drehe mich Richtung Achtern als eine gewaltiger Ruck durch
das Boot geht. Irgendwas hat uns gerammt! Sofort ist Chaos und Panik angesagt. Ich bin weniger nutzlos als
gedacht, binde mir geistesgegenwärtig ein Tau um den Bauch und helfe einem verletzen Matrosen. Das Deck ist
schief. Ich starre über das Deck, aber da wo der Kapitän zuvor stand ist nur Luft. Und aus der Richtung höre ich
Wasser und Schreie. Leute rennen an mir vorbei. Admunson der Norweger, hat sich auch ein Seil um gebunden. Er
läuft und wirft sich über Bord ins Wasser. Es dauert nur wenige Minuten. Gemeinsam ziehen wir den Norweger, der
unseren Kapitän in den Armen hält nach oben. Ein paar Matrosen, haben es geschafft sich an den Seilen fest zu
halten und wir können sie retten. Niemand geht uns an diesem Tag verloren. Aber wir sitzen fest.
\\
\subsection{
  Im Packeis.
}
\noindent Erzählt es nicht der Besatzung, aber ich bin kein Seemann. Aber wenn mir jemand gesagt hätte, dass es sehr
schnell gehen kann, dass man im Eis der Arktis stecken bleiben kann, dann hätte ich gedacht das mich hier jemand
veralbert. Ich dachte es gibt eine riesige Insel oder eine Kontinent komplett aus Eis, oder mit Eis bedeckt und
darum herum schwimmen Eisberge, die nach Lust und Laune die eine oder andere Titanic aufschlitzen. Aber das das
Packeis viel tückischer ist und sich überaschend schnell verschieben kann und so ein Schiff regelrecht in die
Mangel nimmt, dass hätte ich nicht gedacht. Wir sitzen fest. Über die Nacht wurde aus, auf einer Sand bzw.
Eisbank aufgelaufen ein Eisteppich rund um uns herum. Soweit das Auge reicht. So als hätte man das bot mit
mehreren Helikoptern über das Eis gehoben und mitten auf dem Kontinent abgesetzt. Wir hatten trotz der schlechte
Bedingungen trotzdem Glück im Unglück, dass die Männer, die ins Wasser gefallen sind so schnell gerettet werden
konnten. Leichte Erfrierungen. Admunson hat mir heute morgen das Logbuch in die Hand gedrückt, der Kapitän ist
zu erschöpft um seinen Dienst bis auf weiteres zu verüben. Er will, dass ich
Buch führer. Mein Norwegisch ist zu schlecht um abzulehen. Seis drum. Ich beginne zu schreiben. Dann breche ich
direkt wieder ab.
Ich bitte Koch alle Offiziere im Kartenraum zum versammeln. Hier und jetzt erkläre ich ihnen, dass ich nicht der
bin, für den sie mich gehalten haben. Schweigen von den einen, ein leicht enttäuscher Koch, wir sind Freunde
geworden und ich habe sein Vetrauen missbraucht. Admunson ist belustigt. Deswegen sei mein Norwegisch so
beschissen. Ich antworten auf Norwegisch, denn das kann ich mittlerweile okayisch.
Trotz dieser Truthbomb hat keiner wirklich lust mich über Bord zu werfen. Es hilft auch, dass ich mich mit allen
so gut es geht angefreundet habe - ursprünglich um den Kurs der Reise und die vielen Sprachen zu lernen - aber
nun in unserer Situation ist eine Dolmetscher nicht das schlechteste.
\\
\noindent Die Wissenschaftlern übernehmen das Wort. Sie erklären, dass wir vermutlich entweder in den kommenden Tagen
wieder frei kommen, oder hier für Wochen oder sogar Monate fest sitzen werden. Cook weist darauf hin dass wir
nicht in Sauß und Brauß leben können, aber für ein solches Szenario vorgesorgt und ausreichend Vorräte an Bord
haben. Ich nicke, so als wäre ich noch der Expeditionsleiter. Danke allen und bitte Cook und Admunson zu bleiben
während die anderen ihren Aufgaben entlassen sein. Ich muss mir kurz ein Grinsen verkneifen, als alle ohne
Murren wirklich gehen und mir aufmunternd zu nicken. Cook und Admunson blicken mich an. Ich spreche ein Wort
fragend aus. Skorbut?\\
\subsection{
  Der Klub des Pinguins.
}
\noindent 
Das coolste Tier auf der Welt ist vermutlich der gemeine Waldundwiesenhund. Danach das Kaiserbartäffchen. Aber
danach ist bestimmt der Pinguin das coolste Tier auf der Welt. Auf grund dessen behaupte ich jetzt mal, dass die
Belica riesige Mengen unverderblicher frischen Orange und Zitronen an Bord hatte. Pinguine und Robben haben wir
zwar
gefangen aber nur um sie zu streicheln und knuddeln. Eine Robbe nannten wir Feven.
\\
\subsection{
  Heiligabend am Ende der Welt
}
\noindent 
Seit ein paar wochen waren wir nun im Eis gestrandet. An einem Abend legte sich einer der Männer im Schlafsack
auf das Eis neben das Schiff. Er wollte wenigstens eine Nacht weg von dem Gestank des Schiffs und der
Mannschaftskajüte. Ein Anderer, der in dieser Nacht Wache hatte sah die große schwarze Siloutte auf dem Eis. Ein
so großer Seelöwe bedeutet viel Fett und Fleisch. Der Mann lädt sein Gewehr und zielt unwissend auf seinen
Schiffskollegen. Im letzten Moment hält er inne. Er sichert die Waffe und lässt das Tier leben. Es ist
Weihnachten. Das Fest der Liebe und des Friedens.
Auf dem Speißeplan stand alles was die Konservendosen der Vorratsräume zu bieten hatten. Ein hoch auf einen
frischen Apfel. Trotzdem hatten wir uns im Kartenraum ein Weihnachtsfest eingerichtet. Es wurde gelacht und
gegessen. Gestritten und Vergeben und gessoffen, nicht dass der Alkohol schlecht wird.
\\
\subsection{
  Nachwort
}
\noindent 
So liebe Leute. Ihr wisst man soll aufhören wenn es am Schönsten ist. Ich hoffe das nun niemand zu enttäuscht
ist, dass die Geschichte hier an einem Cliffhänger endet. Aber das Weihnachtsfest ist leider für die armen
Teufel, die da in die Antarktis auf gebrochen waren nur eine Etappe in einer noch viel längeren Reise. Und wenn
ich jetzt die ganze Geschichte erzählen müsste, dann würden wir morgen früh noch da sitzten. Es kann auch sein,
dass ich vielleicht das eine oder andere Detail falsch wieder gegeben habe. Aber wichtig ist, wenn ihr die Story
interessant fandet, dann könnt ihr euch schlau machen über die Echte Expedition der Belgica. Ich hoffe dir geht
es gut, und du hast alles zum Weihnachten bekommen, dass du dir gewünscht hast.
\begin{center}Fröhliche Weihnachten\end{center}
  \\

  
\section{English}
  \subsection{Foreword.}
  \noindent 
  Once upon a time, long, long ago. In a far-off place. At an unknown time. So,
  about a few years ago, 1-2 weeks before Christmas. Due to a lack of time, money, and ability in German spelling,
  I decided to honor my loved ones with a piece of work, the newest iteration of which now lies in your hands,
  dear reader. A fixed idea became a tradition, and now I've been sending out the good news for several years. As
  every year, turn off your brain, it's all the purest truth and certainly not made up. Or is it? And above all,
  whoever finds spelling mistakes may count them, square the result, and write it on a napkin in their favorite
  restaurant. Have fun with that. Now that that's out of the way, I'd like to ask everyone to raise their glass
  and yes, keep it up, a little longer, yes, now cheers everyone. And now I wish you much fun reading my nonsense,
  which certainly did not come about under the influence of consciousness-expanding hop tea. (God, the German
  language is something beautiful, the AI will be delighted to translate this into English).\\
  \subsection{A Strange Letter.}
  \noindent 
  As you all know, of course, I am an enthusiastic sailor. Neither the bathtub
  nor the sink is safe from me. Nevertheless, I was slightly surprised when one day a letter found its way into my
  mailbox. Actually, I was hoping for an order of finest chocolate, but alas, it was just a letter. In addition to
  the bitter taste of ink instead of sweet cocoa notes, the sender had also mistaken the person. The letter was
  also in French, and you know I'm an academic - which means I imagine I understand every language a little bit.
  But I shy away from physical work like the devil from holy water. What the message meant, I couldn't outline at
  this point. However, I understood one thing very clearly. First-class train tickets for a journey should not be
  wasted. Even if one's name is neither \textit{Adrien} nor \textit{de Gerlach}. I could deduce from the destination
  of the tickets that the sender was Belgian and not French, because the destination was a Belgian port city. And
  since the chocolate I ordered came from Switzerland and there seems to be quite decent beer in Belgium, I
  democratically decided to accept the invitation. The vote had a total of one participant.\\
  \noindent A week later, I arrived in Antwerp. Incredibly interesting - I had always pronounced it "Antwerben". With a B.
  Funny. Well, since the train tickets were not addressed to me but to another person, I had no hopes of getting
  further than here. My plan was to see the city and the insides of some Belgian beer mugs, and then travel back.
  Sad anecdote on the side - the return flight ticket costs, of course, a fraction of the trains. Did you know
  that
  Antwerp is the second-largest city in Belgium? For lunch, I had mussels with fries, and if I had known what was
  in
  store for me, I might have enjoyed a second beer. Okay, a third beer. But it was a vacation. Don't judge me!\\
  \noindent In good spirits, after the meal and a small sightseeing tour (According to their own claim, Antwerp has the
  oldest museum in the world!), I followed the instructions from the letter, which I still hadn't completely
  translated. My path led me to the \textit{largest port in Europe} (My goodness, what's up with Antwerp?). Past
  huge
  cruise ships, my phone, or rather Google Maps, guided me to a magnificent, albeit not entirely unusual ship.
  Busy
  activity around, on, and in front of the sailing ship led me to suspect that the departure, or as we \textit{sea
    dogs} would say, was imminent. Interested and slightly out of place, I stood there. With my suitcase and
  the
  letter. Belgian sailors were shouting things at each other. Loading and unloading things.\\
  \noindent I was about to turn back when an old-fashioned looking captain with a cap and everything spotted me. I'll spare
  you and myself repeating the wording because I didn't understand him. In any case, he waved me on board, I
  clumsily stumbled across a narrow wooden bridge and now stood before him, Captain \textit{Georges Lecointe}. He
  saluted, then heartily extended his hand to me. I shook it, not wanting to be impolite. And then I pointed at
  his
  uniform and my tourist attire and attempted a joke in French. The captain began to roar with laughter. Whether
  with or at me, I wasn't aware, but I now also doubted his French skills. A second man quickly joined us, eagerly
  grasped my now free hand and shook me. In old-fashioned English, he introduced himself as Frederik Cook, a
  doctor
  from New York. One could immediately tell that his immigration to the States was quite recent. And when I asked
  in
  German whether he preferred to be called Friedrich or Frederik, he was visibly pleased that I \textit{surprisingly
    also} spoke German. At that moment, I felt hot behind the ears. Also? Next, a very tall, very broadly
  muscular man came to me. He growl-mumbled something in pseudo-French. I looked at the doctor. He addressed Mr.
  Admunson and said in broken English that I could also speak \textit{Norwegian}. Admunson's expression brightened,
  and he now growl-mumbled a snippet at me that was surprisingly more understandable than the French before. And I
  don't speak a word of Norwegian. In my head, I tried to recapitulate. These people thought I was a
  French-speaking
  Belgian who also speaks English and Norwegian and German. I had to suppress a grin. Admunson was still looking
  at
  me as if expecting an answer, I raised my hand as if in a gesture to put an end to this farce, but then the
  captain pulled me behind him into the interior of the ship.\\
  \noindent Through a narrow corridor past cabins, he showed me a few rooms, a dining room, a kind of chart room, and
  finally
  opened a door, pushed me in, said goodbye, and I stood in my cabin. Although the ship was impressively large, I
  was still surprised at how luxurious it looked here. Okay, I practically had a bed, a table, and a window or
  porthole. But there was even another door to a small bathroom. I wouldn't have thought that sailors on an
  old-fashioned sailing ship would get such a luxury cabin. Said and bitten my tongue. The uniform, which was
  clearly intended for me, looked nothing like a sailor's. In my head, I panicked, wondering if I even had more
  tinsel on it than the one I thought was the \textit{Captain}. With shaky legs, I put on the disguise. What's the
  punishment for impersonating a high-ranking naval officer on a sailing ship? Keelhauling? I don't remember what
  that is anymore. But I do remember that it's not pleasant. Well, it's all a big misunderstanding. The,
  hopefully,
  Captain practically forced the stuff on me. As I walk back down the narrow corridor and step out through a small
  door into the open, I'm greeted by applause. Great. I raise my arm reassuringly and head towards the
  (FUUUUUUUUUUUUCK I have more tinsel than that dude) Captain. I clearly have more golden swirls and stripes on my
  jacket, and my hat feels twice as big as his. The doctor formally introduces everyone to me again, handing each
  a
  glass of champagne. So this is Doctor Friedrich Koch, Roald Amundsen the 1st Officer, 2 more officers whose
  names
  I instantly forget, and yes, Captain \textit{Georges Lecointe}. I don't know if you knew this, but on ships, it
  seems normal that the Captain isn't the top boss - that would be me, the expedition leader. At this moment, we
  begin to set sail. I look back longingly. If you squint really hard, you can see a Belgian expedition leader
  desperately doing jumping jacks on the pier because he's been forgotten.\\
  \subsection{On the High Seas.}
  \subsubsubsection{My First Week at Sea.}
  \noindent By now, I no longer had to vomit every minute. No, really, you get used to the rocking excellently. It turns
  out
  that with my cabin, I'm actually far ahead of the average sailor. Below deck, there's a stinking, humid warm
  room
  that reminds one of a dormitory in a hostel. Hammocks and sea bags hang in there. The officers and a series of
  scientists have cabins that are slightly larger than my chief suite, but are inhabited by 4 men simultaneously.
  Also, the luxury of a private bath is only enjoyed by me and the Captain - a fine fellow. He chatters at me in
  French non-stop. I'm slowly starting to imagine that I understand a word or two. Just like Norwegian. It's
  actually an incredibly easy language. A Nordic giant comes to you and says something with as few words as
  possible. You nod and maybe add a "yes" in any language. Whatever Amundsen said, he reports it to me only out of
  duty and loyalty to the hierarchy, but he doesn't really need me. He could probably steer the entire ship alone.
  \\
  \noindent Apropos steering. Where are we steering to right now? Well, I have absolutely no idea. I know that there's no
  land in sight in any direction, so we could just as well be in the Mediterranean or the Indian Ocean. Although
  the
  English Channel or maybe even the North Atlantic would make more sense. You can already tell what a great
  navigator I am, right?\\
  \subsubsubsection{My 4th Week at Sea.}
  \noindent I'm slowly beginning to fear that I'm not on a short pleasure trip here. Also, I have a guilty conscience that
  I
  can contribute so little to the journey. Sure, everyone treats me like a king. The dude I'm accidentally
  impersonating is probably the owner of the ship and pays for everything here, but if he's traveling himself,
  doesn't he have a job too? The Captain spends most of his time on deck or in his cabin, so I've now started to
  barricade myself in the chart room for half the night. From the outside, it might look like I'm studying sea
  charts or something, but in reality, I'm reading all the dictionaries and everything that even remotely compares
  German or English to the many other languages on board. It's wonderfully ironic how I'm trying to learn French
  and
  Norwegian on a ship full of Belgians and Norwegians without them noticing that I understand neither language.
  \\
  \noindent During the day, I spend a lot of time on deck listening to people. I often talk with Fritz, the doctor. It's
  pleasant not having to pretend, as he thinks I can speak German very well, so I don't need to come up with a
  complicated story. There have been situations where I could help him write letters in English to his American
  wife. Because I can't directly ask where we're going, I indirectly inquire with the doctor about his supply
  status. If we were to go ashore, then I would know where we are and what our course looks like. For a crazy
  moment, I think I could then approximate our destination. As if I wouldn't flee at the first opportunity on
  land.
  \textit{Unfortunately}, we're so well stocked with supplies that a stop isn't necessary. So I have to keep
  observing
  and listening and pretending I know what's going on.
  \\
  \subsubsubsection{My 5th Week at Sea.}
  \noindent Oh shit!!! Those were my words when I went on deck this morning. Or rather, \textit{Oh, putain de merde}. Yes, my
  French is slowly becoming useful enough to communicate with people. Don't laugh at me, but \textit{Adrien de
    Gerlach}, the person I'm impersonating, may have a French-sounding name, but apparently comes from the
  German-speaking part of the country. And because everyone immediately realized that my French isn't good enough
  to
  buy bread, they just dismissed me as one of those stupid German-Belgians. And no one thought that was weird.
  Back
  to my excellent curse in French. An ice floe is drifting past us on the port side. I think I know where we're
  heading, southward. Waaay southward. The name of the ship has seemed familiar to me all along, but now it hits
  me
  like a ton of bricks. I'm on the Belgica!\\
  \subsubsubsection{My Last Week at Sea.}
  \noindent It was never very warm on board. Or rather, it was warm in my cabin, which has the exhaust pipes of the ovens
  as
  underfloor heating. But right in the corridor next to it or on deck, it became unmistakably clear that you're on
  the high seas. But the temperatures were never extremely cold, just fresh, invigorating. But for a few days now,
  you can tell that we're getting closer and closer to the expedition's destination. By now, I know what's in
  store
  for us. The Belgica is one of those ships that set out to reach the South Pole. I've really gotten myself into
  something here. With a sour expression, I leave my cabin and go on deck. Gone are the days for the fancy
  uniform.
  A long, warm coat has long since been hanging over my shoulders. A thick knit cap adorns my head. I'm in a bad
  mood because, of course, I can't just admit now that I'm not \textit{Adrien de Gerlach}. Consequently, I can't
  simply make the mission turn back, even though we're steering nose-first into the pack ice of Antarctica. And
  every time I cautiously hint whether people are worried about the ice, they only answer with a death-defying
  coolness of seafarers who have already fought against giant krakens, the maelstrom, and the Klabautermann. The
  Captain greets me beaming and we exchange a few words. Then he holds a map in front of my nose and wants to know
  my opinion. Jokingly, I say \textit{Northward}? Howling laughter from the crew at my excellent joke. I just nod
  and
  close my eyes. I count to 10. I open my eyes. Everything is fine. I nod to the Captain and turn towards the
  stern
  when a tremendous jolt goes through the boat. Something has rammed us! Immediately, chaos and panic ensue. I'm
  less useless than I thought, quickly tie a rope around my waist, and help an injured sailor. The deck is
  slanted.
  I stare across the deck, but where the Captain was standing before is only air. And from that direction, I hear
  water and screams. People run past me. Amundsen, the Norwegian, has also tied a rope around himself. He runs and
  throws himself overboard into the water. It only takes a few minutes. Together we pull up the Norwegian, who is
  holding our Captain in his arms. A few sailors have managed to hold onto the ropes, and we can rescue them. We
  don't lose anyone that day. But we're stuck.\\
  \subsection{In the Pack Ice.}
  \noindent Don't tell the crew, but I'm not a sailor. But if someone had told me that it could happen very quickly to get
  stuck in the Arctic ice, I would have thought someone was pulling my leg. I thought there was a huge island or
  continent completely made of ice, or covered with ice, and around it float icebergs that slice open the
  occasional
  Titanic at will. But that the pack ice is much more treacherous and can shift surprisingly quickly, literally
  putting a ship through the wringer, I wouldn't have thought. We're stuck. Overnight, running aground on a
  sandbank
  or ice bank turned into an ice carpet all around us. As far as the eye can see. As if several helicopters had
  lifted the boat over the ice and set it down in the middle of the continent. Despite the bad conditions, we were
  still lucky in our misfortune that the men who fell into the water could be rescued so quickly. Mild frostbite.
  This morning, Amundsen handed me the logbook; the captain is too exhausted to perform his duties for the time
  being. He wants me to keep the book. My Norwegian is too poor to refuse. So be it. I start writing. Then I
  immediately stop again. I ask Cook to gather all officers in the chart room. Here and now I explain to them that
  I'm not who they thought I was. Silence from some, a slightly disappointed Cook - we've become friends and I've
  abused his trust. Amundsen is amused. That's why my Norwegian is so shitty. I answer in Norwegian, because I can
  do that okay-ish by now. Despite this truth bomb, no one really wants to throw me overboard. It also helps that
  I've made friends with everyone as best I could - originally to learn the course of the journey and the many
  languages - but now in our situation, an interpreter is not the worst thing to have.\\
  \noindent The scientists take over. They explain that we will probably either be free again in the coming days, or be
  stuck
  here for weeks or even months. Cook points out that we can't live in luxury, but we have prepared for such a
  scenario and have sufficient supplies on board. I nod, as if I were still the expedition leader. Thank everyone
  and ask Cook and Amundsen to stay while the others are dismissed to their duties. I have to suppress a grin for
  a
  moment when everyone really leaves without grumbling and nods encouragingly at me. Cook and Amundsen look at me.
  I
  questioningly utter one word. Scurvy?\\
  \subsection{The Penguin Club.}
  \noindent The coolest animal in the world is probably the common forest and meadow dog. After that, the emperor tamarin.
  But after that, the penguin is surely the coolest animal in the world. Based on that, I'm going to claim now
  that
  the Belgica had huge amounts of non-perishable fresh oranges and lemons on board. We did catch penguins and
  seals,
  but only to pet and cuddle them. We named one seal Feven.\\
  \subsection{Christmas Eve at the End of the World}
  \noindent We had been stranded in the ice for a few weeks now. One evening, one of the men laid down in a sleeping bag on
  the ice next to the ship. He wanted to spend at least one night away from the stench of the ship and the crew's
  cabin. Another, who was on watch that night, saw the large black silhouette on the ice. Such a large sea lion
  means a lot of fat and meat. The man loads his gun and unknowingly aims at his shipmate. At the last moment, he
  pauses. He secures the weapon and lets the animal live. It's Christmas. The feast of love and peace. On the menu
  was everything the canned goods of the storage rooms had to offer. Here's to a fresh apple. Nevertheless, we had
  set up a Christmas celebration in the chart room. There was laughter and eating. Arguing and forgiving and
  drinking, not that the alcohol goes bad.\\
  \subsection{Afterword}
  \noindent So, dear people. You know you should stop when it's at its best. I hope now no one is too disappointed that the
  story ends here on a cliffhanger. But the Christmas celebration is unfortunately only one stage in an even
  longer
  journey for the poor devils who had set out for Antarctica. And if I had to tell the whole story now, we'd still
  be sitting here tomorrow morning. It's also possible that I might have gotten one or another detail wrong. But
  what's important is, if you found the story interesting, you can learn about the real expedition of the Belgica.
  I
  hope you're doing well, and you got everything for Christmas that you wished for.\\
  \begin{center}Merry Christmas\end{center}

\section{Español}
    \subsection{Prólogo.}
    \noindent 
    Érase una vez, hace mucho, mucho tiempo. En un lugar muy lejano. En una época
    desconocida. Es decir, hace unos años, 1-2 semanas antes de Navidad. Por falta de tiempo, dinero y habilidad en
    la
    ortografía alemana, decidí obsequiar a mis seres queridos con una obra cuya última iteración ahora también está
    en
    tus manos, estimado lector. Una idea fija se convirtió en una tradición y ahora llevo varios años enviando las
    buenas nuevas. Como cada año, apaguen sus cerebros, todo es la pura verdad y ciertamente no inventado. ¿O sí? Y
    sobre todo, quien encuentre faltas de ortografía, que las cuente, las eleve al cuadrado y escriba el resultado
    en
    una servilleta en su restaurante favorito. Que se diviertan. Después de aclarar esto, me gustaría pedir a todos
    que levanten sus copas y sí, manténganlas arriba, un poco más, sí, ahora brindemos todos. Y ahora les deseo que
    disfruten leyendo mis tonterías, que ciertamente no surgieron bajo la influencia de té de lúpulo que expande la
    conciencia. (Dios, el idioma alemán es algo hermoso, la IA se alegrará de traducir esto al inglés).\\
    \subsection{Una carta extraña.}
    \noindent 
    Como todos saben, por supuesto, soy un entusiasta marinero. Ni la bañera ni %
    el
    lavabo están a salvo de mí. Sin embargo, me sorprendí ligeramente cuando un día encontré una carta en mi buzón.
    En
    realidad, esperaba un pedido del mejor chocolate, pero desafortunadamente no, era solo una carta. Además del
    sabor
    amargo de la tinta en lugar de las dulces notas de cacao, el remitente también se había equivocado de persona.
    La
    carta estaba además en francés y ustedes saben que soy un académico - lo que significa que me imagino que
    entiendo
    un poco cada idioma. Pero me aparto del trabajo físico como el diablo del agua bendita. No pude descifrar lo que
    significaba el mensaje en ese momento. Sin embargo, entendí una cosa muy claramente. No se deben desperdiciar
    los
    billetes de tren de primera clase para un viaje. Aunque uno no se llame ni \textit{Adrien} ni \textit{de Gerlach}.
    Pude deducir del destino de los billetes que el remitente era belga y no francés, porque el destino era una
    ciudad
    portuaria belga. Y como el chocolate que pedí venía de Suiza y parece que en Bélgica hay cerveza bastante
    decente,
    decidí democráticamente aceptar la invitación. La votación tuvo un total de un participante.\\
    \noindent Una semana después llegué a Amberes. Súper interesante - siempre lo había pronunciado "Antwerben". Con una B.
    Gracioso. Bueno, como los billetes de tren no estaban dirigidos a mí sino a otra persona, no tenía esperanzas de
    llegar más allá de aquí. Mi plan era ver la ciudad y el interior de algunas jarras de cerveza belga y luego
    volver. Anécdota triste al margen - el billete de avión de vuelta cuesta, por supuesto, una fracción de los
    trenes. ¿Sabían que Amberes es la segunda ciudad más grande de Bélgica? Para el almuerzo hubo mejillones con
    papas
    fritas y si hubiera sabido lo que me esperaba, tal vez me habría tomado una segunda cerveza. Bueno, está bien,
    una
    tercera cerveza. Pero eran vacaciones. ¡No me juzguen! De buen humor, después de comer y de un pequeño tour
    turístico (¡En Amberes está, según dicen, el museo más antiguo del mundo!), seguí las indicaciones de la carta,
    que aún no había traducido completamente. Mi camino me llevó al \textit{puerto más grande de Europa} (¡Dios mío,
    qué
    pasa con Amberes?). Pasando enormes cruceros, mi teléfono, o mejor dicho Google Maps, me guió hasta un barco
    magnífico, aunque no del todo inusual. El ajetreo alrededor, sobre y frente al velero me hizo suponer que la
    partida, como diríamos nosotros los \textit{lobos de mar}, estaba a punto de ocurrir. Interesado y ligeramente
    fuera
    de lugar, me quedé allí. Con mi maleta y la carta. Los marineros belgas se gritaban cosas. Cargando y
    descargando
    cosas.\\
    \noindent Estaba a punto de dar media vuelta cuando un capitán de aspecto anticuado con gorra y todo me avistó. Les
    ahorro
    a ustedes y a mí repetir las palabras porque no las entendí. En cualquier caso, me hizo señas para que subiera a
    bordo, tropecé torpemente por un estrecho puente de madera y ahora estaba frente a él, el Capitán \textit{Georges
      Lecointe}. Me saludó, luego me tendió cordialmente la mano. La estreché, uno no quiere ser descortés. Y
    luego
    señalé su uniforme y mi atuendo de turista e intenté hacer una broma en francés. El capitán comenzó a reír a
    carcajadas. Si conmigo o de mí, no lo sabía, pero ahora también dudaba de sus conocimientos de francés. Un
    segundo
    hombre se unió rápidamente a nosotros, agarró entusiasmado mi mano ahora libre y me sacudió. En un inglés
    anticuado se presentó como Frederik Cook, médico de Nueva York. Se le notaba directamente que su inmigración a
    los
    Estados Unidos era bastante reciente. Y cuando pregunté en alemán si prefería que lo llamaran Friedrich o
    Frederik, se alegró visiblemente de que yo \textit{sorprendentemente también} hablara alemán. En ese momento sentí
    calor detrás de las orejas. ¿También? A continuación, se me acercó un hombre muy alto y muy musculoso. Gruñó
    algo
    en pseudo-francés. Miré al doctor. Este se dirigió al Sr. Amundsen y dijo en un inglés entrecortado que yo
    también
    hablaba \textit{noruego}. El rostro de Amundsen se iluminó y ahora me gruñó un fragmento que sorprendentemente era
    más comprensible que el francés anterior. Y yo no hablo ni una palabra de noruego. En mi cabeza intenté
    recapitular. Esta gente pensaba que yo era un belga francófono que también hablaba inglés, noruego y alemán.
    Tuve
    que contener una sonrisa. Amundsen seguía mirándome como si esperara una respuesta, levanté la mano como en un
    gesto con el que quería poner fin a esta farsa, pero entonces el capitán me arrastró tras él al interior del
    barco.\\
    \noindent A través de un estrecho pasillo pasando por camarotes, me mostró algunas habitaciones, un comedor, una especie
    de
    sala de mapas y finalmente abrió una puerta, me empujó dentro, se despidió y me encontré en mi camarote. Aunque
    el
    barco era impresionantemente grande, me sorprendió lo lujoso que se veía aquí. Bueno, tenía prácticamente una
    cama, una mesa y una ventana o un ojo de buey. Pero incluso había otra puerta que daba a un pequeño baño. No
    habría pensado que los marineros en un velero anticuado recibirían una cabina tan lujosa. Dicho y mordido la
    lengua. El uniforme, que claramente estaba destinado a mí, no se parecía en nada al de un marinero. En mi
    cabeza,
    me pregunté con pánico si incluso tenía más oropel que el que pensaba que era el \textit{Capitán}. Con piernas
    temblorosas, me puse el disfraz. ¿Cuál es el castigo por imitar a un alto dignatario en el mar en un velero?
    ¿Pasar por debajo de la quilla? Ya no recuerdo lo que es eso. Pero recuerdo que no es agradable. Bueno, todo es
    un
    gran malentendido. El, espero, Capitán prácticamente me impuso estas cosas. Cuando vuelvo por el estrecho
    pasillo
    y salgo por una pequeña puerta al aire libre, me recibe un aplauso. Genial. Levanto el brazo para calmar y me
    dirijo hacia el (MIEEEEEERDA tengo más oropel que ese tipo) Capitán. Claramente tengo más espirales y rayas
    doradas en mi chaqueta y mi gorra se siente el doble de grande que la suya. El doctor me presenta formalmente a
    todos de nuevo, y le da a cada uno una copa de champán. Así que este es el Doctor Friedrich Koch, Roald Amundsen
    el primer oficial, otros 2 oficiales cuyos nombres olvido al instante y sí, el Capitán \textit{Georges Lecointe}.
    No
    sé si lo sabían, pero en los barcos aparentemente es normal que el Capitán no sea el jefe supremo - ese soy yo,
    el
    líder de la expedición. En este momento comenzamos a zarpar. Miro hacia atrás con nostalgia. Si uno entrecierra
    los ojos con fuerza, puede ver a un líder de expedición belga haciendo saltos desesperados en el muelle porque
    lo
    han olvidado.\\
    \subsection{En alta mar.}
    \subsubsubsection{Mi primera semana en el mar.}
    \noindent Ya no tenía que vomitar cada minuto. No, en serio, uno se acostumbra excelentemente al balanceo. Resulta que
    con
    mi camarote estoy muy por delante del marinero promedio. Bajo cubierta hay una habitación maloliente, húmeda y
    cálida que recuerda a un dormitorio en un albergue. Allí cuelgan hamacas y bolsas de marinero. Los oficiales y
    una
    serie de científicos tienen camarotes que son un poco más grandes que mi suite de jefe, pero están ocupados por
    4
    hombres al mismo tiempo. Además de mí, solo el Capitán tiene el lujo de un baño privado - un buen tipo. Me habla
    sin parar en francés. Poco a poco me imagino que entiendo una que otra palabra. Al igual que el noruego. En
    realidad es un idioma muy fácil. Un gigante nórdico viene a ti y dice algo con la menor cantidad de palabras
    posible. Asientes y tal vez añades un sí en cualquier idioma. Sea lo que sea que Amundsen haya dicho, me lo
    informa solo por deber y lealtad a la jerarquía, pero no me necesita. Probablemente podría dirigir todo el barco
    solo.\\
    \noindent A propósito de dirigir. ¿Hacia dónde nos dirigimos ahora? Bueno, no tengo absolutamente ni idea. Sé que no hay
    tierra a la vista en ninguna dirección y por lo tanto podríamos estar igualmente en el Mediterráneo o en el
    Océano
    Índico. Aunque el Canal de la Mancha o tal vez incluso el Atlántico Norte tendría más sentido. ¿Ya se dan cuenta
    de qué navegante tan increíble soy, verdad?\\
    \subsubsubsection{Mi cuarta semana en el mar.}
    \noindent Lentamente temo que no estoy aquí en un corto viaje de placer. Además, tengo remordimientos de que puedo
    contribuir tan poco al viaje. Claro, todos me tratan como a un rey. El tipo al que represento accidentalmente es
    probablemente el dueño del barco y paga todo aquí, pero si él mismo viaja, también debe tener un trabajo, ¿no?
    El
    Capitán pasa la mayor parte del tiempo en cubierta o en su camarote, así que ahora he empezado a atrincherarme
    en
    la sala de mapas durante la mitad de la noche. Desde fuera puede parecer que estoy estudiando cartas náuticas o
    algo así, pero en realidad estoy leyendo todos los diccionarios y todo lo que remotamente compara el alemán o el
    inglés con los muchos otros idiomas a bordo. Es maravillosamente irónico cómo intento aprender francés y noruego
    en un barco lleno de belgas y noruegos sin que se den cuenta de que no entiendo ni uno ni otro idioma.\\
    \noindent Durante el día, paso mucho tiempo en cubierta escuchando a la gente. A menudo también hablo con Fritz, el
    médico.
    Es agradable no tener que fingir, ya que él piensa que puedo hablar alemán muy bien y así no necesito inventar
    una
    historia complicada. Ya ha habido situaciones en las que pude ayudarlo a escribir cartas en inglés a su esposa
    americana. Como no puedo preguntar directamente a dónde vamos, indago indirectamente con el médico sobre su
    estado
    de suministros. Si fuéramos a tierra, sabría dónde estamos y cómo es nuestro rumbo. Por un momento loco, pienso
    que entonces podría aproximar nuestro destino. Como si no huyera en la primera oportunidad en tierra.
    \textit{Desafortunadamente}, estamos tan bien abastecidos que no es necesario detenerse. Así que tengo que seguir
    observando y escuchando y fingir que sé lo que está pasando.
    \\
    \subsubsubsection{Mi quinta semana en el mar.}
    \noindent ¡Mierda!!! Esas fueron mis palabras cuando subí a cubierta esta mañana. O más bien \textit{Oh, putain de merde}.
    Sí, mi francés es lo suficientemente útil como para comunicarme con la gente. No se rían de mí, pero \textit{Adrien
      de
      Gerlach}, la persona que estoy personificando, tiene un nombre que suena francés, pero aparentemente viene
    de
    la parte germanoparlante del país. Y como todos se dieron cuenta inmediatamente de que mi francés no es
    suficiente
    ni para comprar pan, me han descartado como uno de esos estúpidos belgas alemanes. Y a nadie se le ocurrió que
    eso
    era extraño. Volviendo a mi excelente maldición en francés. Un témpano de hielo pasa por nuestro lado de babor.
    Creo que sé hacia dónde nos dirigimos, hacia el sur. Muuuy hacia el sur. El nombre del barco me ha parecido
    familiar todo el tiempo, pero ahora me cae como escamas de los ojos. ¡Estoy en el Belgica!\\
    \subsubsubsection{Mi última semana en el mar.}
    \noindent Nunca hizo mucho calor a bordo. O mejor dicho, sí hacía calor en mi camarote, que tiene los tubos de escape de
    los hornos como calefacción por suelo radiante. Pero justo en el pasillo de al lado o en cubierta, quedaba
    inequívocamente claro que uno se encuentra en alta mar. Sin embargo, las temperaturas nunca fueron
    extremadamente
    frías, sino frescas, vigorizantes. Pero desde hace unos días se nota que nos acercamos cada vez más al objetivo
    de
    la expedición. Ya sé lo que nos espera. El Belgica es uno de esos barcos que zarparon para alcanzar el Polo Sur.
    Me he metido en un buen lío. Con cara agria, salgo de mi camarote y subo a cubierta. Se acabaron los tiempos del
    elegante uniforme. Hace tiempo que llevo un abrigo largo y cálido sobre los hombros. Un grueso gorro de punto
    adorna mi cabeza. Estoy de mal humor porque, por supuesto, no puedo admitir ahora que no soy \textit{Adrien de
      Gerlach}. En consecuencia, no puedo simplemente hacer que la misión dé la vuelta, aunque nos estemos
    dirigiendo de narices hacia el hielo del Antártico. Y cada vez que insinúo con cautela si la gente está
    preocupada
    por el hielo, responden solo con una frialdad desafiante de marineros que ya han luchado contra krakens
    gigantes,
    el maelstrom y el Klabautermann. El capitán me saluda radiante y intercambiamos unas palabras. Luego me pone un
    mapa delante de la nariz y quiere saber mi opinión. En broma digo \textit{¿Hacia el norte?} Risas estridentes de
    la
    tripulación por mi excelente broma. Solo asiento y cierro los ojos. Cuento hasta 10. Abro los ojos. Todo está
    bien. Asiento al capitán y me giro hacia popa cuando una sacudida violenta recorre el barco. ¡Algo nos ha
    embestido! De inmediato, cunde el caos y el pánico. Soy menos inútil de lo que pensaba, me ato rápidamente una
    cuerda alrededor de la cintura y ayudo a un marinero herido. La cubierta está inclinada. Miro fijamente la
    cubierta, pero donde antes estaba el capitán solo hay aire. Y de esa dirección oigo agua y gritos. La gente
    corre
    a mi lado. Amundsen, el noruego, también se ha atado una cuerda. Corre y se lanza por la borda al agua. Solo
    tarda
    unos minutos. Juntos izamos al noruego, que sostiene a nuestro capitán en sus brazos. Algunos marineros han
    logrado aferrarse a las cuerdas y podemos rescatarlos. No perdemos a nadie ese día. Pero estamos atrapados.\\
    \subsection{En el hielo.}
    \noindent No se lo digan a la tripulación, pero no soy marinero. Pero si alguien me hubiera dicho que uno puede quedar
    atrapado en el hielo del Ártico muy rápidamente, habría pensado que alguien me estaba tomando el pelo. Pensé que
    había una isla enorme o un continente completamente de hielo, o cubierto de hielo, y alrededor nadaban icebergs
    que rajaban uno que otro Titanic a su antojo. Pero que el hielo marino es mucho más traicionero y puede moverse
    sorprendentemente rápido, atrapando literalmente un barco, eso no lo habría pensado. Estamos atrapados. Durante
    la
    noche, lo que era encallar en un banco de arena o hielo se convirtió en una alfombra de hielo a nuestro
    alrededor.
    Hasta donde alcanza la vista. Como si hubieran levantado el barco con varios helicópteros sobre el hielo y lo
    hubieran dejado en medio del continente. A pesar de las malas condiciones, tuvimos suerte dentro de la desgracia
    de que los hombres que cayeron al agua pudieran ser rescatados tan rápidamente. Leves congelaciones. Amundsen me
    ha entregado el libro de bitácora esta mañana, el capitán está demasiado agotado para cumplir con su deber por
    el
    momento. Quiere que yo lleve el libro. Mi noruego es demasiado malo para rechazar. En fin. Empiezo a escribir.
    Luego paro de inmediato. Pido a Cook que reúna a todos los oficiales en la sala de mapas. Aquí y ahora les
    explico
    que no soy quien creían que era. Silencio de algunos, un Cook ligeramente decepcionado, nos hemos hecho amigos y
    he abusado de su confianza. Amundsen está divertido. Por eso mi noruego es tan malo. Respondo en noruego, porque
    ahora puedo hacerlo más o menos bien. A pesar de esta bomba de la verdad, nadie tiene realmente ganas de tirarme
    por la borda. También ayuda que me haya hecho amigo de todos lo mejor que pude - originalmente para aprender el
    rumbo del viaje y los muchos idiomas - pero ahora en nuestra situación, un intérprete no es lo peor.\\
    \noindent Los científicos toman la palabra. Explican que probablemente o bien nos liberaremos en los próximos días, o
    bien
    nos quedaremos atrapados aquí durante semanas o incluso meses. Cook señala que no podemos vivir a cuerpo de rey,
    pero que hemos previsto tal escenario y tenemos suficientes provisiones a bordo. Asiento, como si todavía fuera
    el
    líder de la expedición. Agradezco a todos y pido a Cook y Amundsen que se queden mientras los demás son
    liberados
    de sus tareas. Tengo que contener una sonrisa cuando todos se van realmente sin quejarse y me asienten con
    ánimo.
    Cook y Amundsen me miran. Pronuncio una palabra interrogativamente. ¿Escorbuto?\\
    \subsection{El Club del Pingüino.}
    \noindent Probablemente el animal más genial del mundo sea el perro común de bosque y pradera. Después el tití emperador.
    Pero después de eso, seguramente el pingüino sea el animal más genial del mundo. Basándome en eso, ahora afirmo
    que el Belgica tenía a bordo enormes cantidades de naranjas y limones frescos no perecederos. Aunque capturamos
    pingüinos y focas, fue solo para acariciarlos y abrazarlos. A una foca la llamamos Feven.\\
    \subsection{Nochebuena en el fin del mundo}
    \noindent Llevábamos unas semanas atrapados en el hielo. Una noche, uno de los hombres se acostó en un saco de dormir
    sobre
    el hielo junto al barco. Quería pasar al menos una noche lejos del hedor del barco y del camarote de la
    tripulación. Otro, que estaba de guardia esa noche, vio la gran silueta negra sobre el hielo. Un león marino tan
    grande significa mucha grasa y carne. El hombre carga su rifle y apunta sin saberlo a su compañero de barco. En
    el
    último momento se detiene. Asegura el arma y deja vivir al animal. Es Navidad. La fiesta del amor y la paz. En
    el
    menú estaba todo lo que las latas de conserva de los almacenes tenían para ofrecer. Un brindis por una manzana
    fresca. A pesar de todo, habíamos organizado una fiesta de Navidad en la sala de mapas. Hubo risas y comida.
    Discusiones y perdones, y bebida, no sea que el alcohol se eche a perder.\\
    \subsection{Epílogo}       %
    \noindent Bueno, querida gente. Ya saben que hay que parar cuando mejor se está. Espero que nadie esté demasiado
    decepcionado porque la historia termine aquí en un cliffhanger. Pero la fiesta de Navidad es, desafortunadamente
    para los pobres diablos que se habían embarcado hacia la Antártida, solo una etapa en un viaje mucho más largo.
    Y
    si ahora tuviera que contar toda la historia, estaríamos sentados aquí hasta mañana por la mañana. También puede
    ser que haya reproducido incorrectamente algún que otro detalle. Pero lo importante es que, si les pareció
    interesante la historia, pueden informarse sobre la verdadera expedición del Belgica. Espero que estés bien y
    que
    hayas recibido todo lo que deseabas para Navidad.\\
    \begin{center}Feliz Navidad\end{center}


\end{document}                  %
